\begin{center}
	\textbf{{\Huge ĐỀ LUYỆN THI TOÁN QUỐC TẾ AMC 8}}
\end{center}
\subsubsection*{Phần I. Câu hỏi trắc nghiệm (Thí sinh chọn một phương án đúng)}
\setcounter{bt}{0}\setlist[enumerate,1]{label=\quad\bfseries\Alph*.}
\begin{baitap}
Tính giá trị của biểu thức:
$$P = 20 - (2 \times 4 + 6 : 2)$$
\begin{enumerate}
	\item $9$.
	\item $11$.
	\item $13$.
	\item $15$.
	\item $17$.
\end{enumerate}
\begin{traloi}
	\dapan{A}
	Thực hiện theo thứ tự phép tính trong ngoặc trước: $2 \times 4 = 8$ và $6 : 2 = 3$, do đó trong ngoặc là $8 + 3 = 11$. Khi đó $P = 20 - 11 = 9$.
\end{traloi}
\end{baitap}

\begin{baitap}
Trung bình cộng của năm số nguyên liên tiếp bằng $15$. Số lớn nhất trong năm số đó là:
\begin{enumerate}
	\item $15$.
	\item $16$.
	\item $17$.
	\item $18$.
	\item $19$.
\end{enumerate}
\begin{traloi}
	\dapan{C}
	Gọi năm số nguyên liên tiếp là $n-2, n-1, n, n+1, n+2$. Trung bình cộng của chúng bằng $n = 15$. Số lớn nhất trong năm số đó là $n + 2 = 15 + 2 = 17$.
\end{traloi}
\end{baitap}

\begin{baitap}
Một hình chữ nhật có chu vi bằng $36\text{ cm}$. Nếu chiều dài giảm đi $2\text{ cm}$ và chiều rộng tăng thêm $2\text{ cm}$ thì hình chữ nhật trở thành một hình vuông. Diện tích ban đầu của hình chữ nhật là:
\begin{enumerate}
	\item $72\text{ cm}^2$.
	\item $77\text{ cm}^2$.
	\item $80\text{ cm}^2$.
	\item $81\text{ cm}^2$.
	\item $85\text{ cm}^2$.
\end{enumerate}
\begin{traloi}
	\dapan{B}
	Nửa chu vi của hình chữ nhật là $36 : 2 = 18\text{ cm}$. Khi chiều dài bớt $2\text{ cm}$ và chiều rộng thêm $2\text{ cm}$ thì được hình vuông có cạnh là $18 : 2 = 9\text{ cm}$. Chiều dài ban đầu là $9 + 2 = 11\text{ cm}$, chiều rộng ban đầu là $9 - 2 = 7\text{ cm}$. Diện tích ban đầu là $11 \times 7 = 77\text{ cm}^2$.
\end{traloi}
\end{baitap}

\begin{baitap}
Một người đi xe đạp từ A đến B với vận tốc $12\text{ km/h}$, sau đó quay trở về từ B về A theo cùng con đường với vận tốc $20\text{ km/h}$. Vận tốc trung bình của cả chuyến đi và về là:
\begin{enumerate}
	\item $15\text{ km/h}$.
	\item $16\text{ km/h}$.
	\item $14{,}5\text{ km/h}$.
	\item $15{,}5\text{ km/h}$.
	\item $16{,}5\text{ km/h}$.
\end{enumerate}
\begin{traloi}
	\dapan{A}
	Gọi quãng đường AB là $d\text{ (km)}$. Thời gian đi là $\frac{d}{12}\text{ (h)}$, thời gian về là $\frac{d}{20}\text{ (h)}$. Tổng thời gian cả đi và về là $\frac{d}{12} + \frac{d}{20} = \frac{2d}{15}\text{ (h)}$. Vận tốc trung bình là $v_{tb} = \frac{2d}{\frac{2d}{15}} = 15\text{ km/h}$.
\end{traloi}
\end{baitap}

\begin{baitap}
Có bao nhiêu số tự nhiên có ba chữ số đôi một khác nhau được lập từ các chữ số thuộc tập hợp $\{1; 2; 3; 4; 5\}$ sao cho số đó chia hết cho $5$?
\begin{enumerate}
	\item $12$.
	\item $16$.
	\item $20$.
	\item $24$.
	\item $60$.
\end{enumerate}
\begin{traloi}
	\dapan{A}
	Để số chia hết cho $5$, chữ số hàng đơn vị bắt buộc phải là $5$ (có $1$ cách chọn). Do các chữ số đôi một khác nhau: chữ số hàng trăm có $4$ cách chọn từ $\{1; 2; 3; 4\}$, chữ số hàng chục có $3$ cách chọn từ $3$ chữ số còn lại. Số các số thỏa mãn là $4 \times 3 \times 1 = 12$ số.
\end{traloi}
\end{baitap}
